\documentclass[a4paper,14pt]{extarticle}
\usepackage{styles}

\begin{document}

This subsection compares the Section~3 benchmarks with results from the library. Both run paired GPU-vs-CPU experiments with the same seeds and configurations. The goal is to check whether the library reproduces the original results.

\subsubsection*{Seed Reproducibility and Environment} Both implementations use the same seeding: \texttt{seed = i} for $i = 0, 1, \dots, n{-}1$. In the Sec.~3 code this loop is in \texttt{statistical\_analysis\_base.py}; in the library it is in \texttt{statistics/experiment.py}. Each seed is passed to both GPU and CPU runs, so every pair starts from the same initial population.

For TSP, the city layout comes from a fixed graph seed, so all runs use the same distance matrix. The bootstrap CI also uses a fixed seed across 10{,}000 resamples. Any difference in reported results therefore comes from the implementation, not randomness.

All experiments were executed on the same machine: an NVIDIA GeForce RTX 2050 laptop GPU (4.0\,GB, Compute Capability 8.6), 12 CPU cores, 16\,GB RAM, running Linux 6.6 under WSL2 with Python 3.12.

\subsubsection*{JIT Warm-Up and Discarded Runs}

Numba's CUDA JIT compiler translates Python kernel code into GPU machine code on first invocation. In the Sec.~3 benchmarks, this compilation occurs during the first measured run (seed~0), producing a wall time spike of 2--3$\times$ the steady-state value. The effect is particularly pronounced in the library, which compiles six separate kernel modules (\texttt{GAMemory\-Manager}, \texttt{Fitness\-Evaluator}, \texttt{Selection\-Kernels}, \texttt{Crossover\-Kernels}, \texttt{Mutation\-Kernels}, \texttt{Population\-Kernels}) compared to four inline \texttt{@cuda.jit} functions in the Sec.~3 code. For example, the library's XOR run~0 took 1{,}166\,ms vs.\ a steady-state mean of 402\,ms (2.9$\times$ spike), inflating the overall standard deviation by a factor of three.

To eliminate this bias, the library executes one discarded GPU+CPU pair with a dedicated warm-up seed before the measured runs begin. This is standard practice for GPU benchmarks - it ensures that all kernel compilation occurs outside the measurement window. The Sec.~3 benchmarks do not apply this correction, which explains why their reported standard deviations are higher and confidence intervals wider. The impact of this strategy is quantified in the benchmark sections below.

\subsubsection*{XOR Perceptron Problem Benchmarking}

Both frameworks were configured identically for the benchmark: population size 100, 500 generations, mutation rate 0.1, 30 paired runs (seeds 0--29). Table~\ref{tab:xor_comparison} summarizes the statistical results.

\begin{table}[H]
\centering
\caption{XOR perceptron benchmark comparison (30 runs, pop=100, gen=500)}
\label{tab:xor_comparison}
\begin{tabular}{lcc}
\toprule
\textbf{Metric} & \textbf{Sec.\,3 Benchmarks} & \textbf{Library} \\
\midrule
GPU mean wall time (ms) & 368.3 $\pm$ 44.5 & 319.7 $\pm$ 24.7 \\
CPU mean wall time (ms) & 555.2 $\pm$ 59.9 & 444.1 $\pm$ 13.7 \\
Speedup & 1.51$\times$ & 1.39$\times$ \\
Speedup 95\% CI & [0.93$\times$, 1.51$\times$] & [1.35$\times$, 1.43$\times$] \\
Time $p$-value & $< 0.0001$ & $< 0.0001$ \\
Transfer overhead & 29.2\% & 31.5\% \\
GPU evals/sec & 201{,}153 & 236{,}755 \\
Fitness (both) & 2.0 $\pm$ 0.0 & 2.0 $\pm$ 0.0 \\
\bottomrule
\end{tabular}
\end{table}

Both implementations converge to the theoretical XOR fitness ceiling of 2 out of 4 patterns with zero variance, confirming correctness. GPU speedup is significant in both cases ($p < 0.0001$), at 1.51$\times$ and 1.39$\times$ respectively, with overlapping 95\% bootstrap confidence intervals.

Transfer overhead is similar (29.2\% vs.\ 31.5\%), showing that pinned-memory transfers work the same way in both. The library reports higher GPU throughput (236{,}755 vs.\ 201{,}153 evals/sec) because the warm-up run removes JIT compilation from the measurements. This also shows in the standard deviation: 24.7\,ms vs.\ 44.5\,ms, giving tighter confidence intervals ([1.35$\times$, 1.43$\times$] vs.\ [0.93$\times$, 1.51$\times$]).

\begin{figure}[H]
\centering
\begin{subfigure}{0.49\textwidth}
\includegraphics[width=\textwidth]{diagrams/sec4/comparison_wall_time_xor.png}
\caption{Wall time comparison}
\label{fig:xor_wall_time}
\end{subfigure}
\hfill
\begin{subfigure}{0.49\textwidth}
\includegraphics[width=\textwidth]{diagrams/sec4/comparison_gpu_breakdown_xor.png}
\caption{GPU time breakdown: compute vs.\ transfer}
\label{fig:xor_gpu_breakdown}
\end{subfigure}
\end{figure}

Figure~\ref{fig:xor_wall_time} shows GPU and CPU wall times in the same range across frameworks, with the library slightly lower due to JIT warm-up. Figure~\ref{fig:xor_gpu_breakdown} shows the same compute/transfer split in both cases (transfer around 29--32\%).

\begin{figure}[H]
\centering
\begin{subfigure}{0.49\textwidth}
\includegraphics[width=\textwidth]{diagrams/sec4/comparison_throughput_xor.png}
\caption{Fitness evaluation throughput}
\label{fig:xor_throughput}
\end{subfigure}
\hfill
\begin{subfigure}{0.49\textwidth}
\includegraphics[width=\textwidth]{diagrams/sec4/comparison_run_by_run_xor.png}
\caption{Run-by-run wall times (30 paired runs)}
\label{fig:xor_run_by_run}
\end{subfigure}
\end{figure}

Figure~\ref{fig:xor_throughput} shows higher GPU throughput in the library (237K vs.\ 201K evals/sec) because JIT overhead is excluded. Figure~\ref{fig:xor_run_by_run} shows stable per-run times in both frameworks, with a tighter band in the library.

\subsubsection*{TSP Problem Benchmarking}

The TSP benchmark used identical configurations across both frameworks: 50 cities, population size 200, 100 generations, tournament size 5, elite rate 0.5, graph seed 42, 50 paired runs (seeds 0--49). Table~\ref{tab:tsp_comparison} presents the results.

\begin{table}[H]
\centering
\caption{TSP benchmark comparison (50 runs, 50 cities, pop=200, gen=100)}
\label{tab:tsp_comparison}
\begin{tabular}{lcc}
\toprule
\textbf{Metric} & \textbf{Sec.\,3 Benchmarks} & \textbf{Library} \\
\midrule
GPU mean wall time (ms) & 496.8 $\pm$ 184.8 & 446.5 $\pm$ 75.3 \\
CPU mean wall time (ms) & 664.2 $\pm$ 61.9 & 504.5 $\pm$ 14.2 \\
Speedup & 1.34$\times$ & 1.13$\times$ \\
Speedup 95\% CI & [1.19$\times$, 1.46$\times$] & [1.08$\times$, 1.19$\times$] \\
Time $p$-value & $< 0.0001$ & $< 0.0001$ \\
Transfer overhead & 23.8\% & 25.1\% \\
GPU evals/sec & 186{,}052 & 157{,}036 \\
\midrule
GPU mean path cost & 9{,}469.9 $\pm$ 604.9 & 9{,}277.6 $\pm$ 634.5 \\
CPU mean path cost & 9{,}889.1 $\pm$ 832.7 & 9{,}713.3 $\pm$ 664.1 \\
Fitness $p$-value & 0.0072 & 0.0072 \\
\bottomrule
\end{tabular}
\end{table}

Both show significant GPU wall-time improvements ($p < 0.0001$). Transfer overhead is nearly identical (23.8\% vs.\ 25.1\%), meaning the CUDA execution path works the same way. Both also show a significant GPU advantage in fitness ($p = 0.0072$), with GPU tours achieving lower mean path cost.

\begin{figure}[H]
\centering
\begin{subfigure}{0.49\textwidth}
\includegraphics[width=\textwidth]{diagrams/sec4/comparison_wall_time_tsp.png}
\caption{Wall time comparison}
\label{fig:tsp_wall_time}
\end{subfigure}
\hfill
\begin{subfigure}{0.49\textwidth}
\includegraphics[width=\textwidth]{diagrams/sec4/comparison_gpu_breakdown_tsp.png}
\caption{GPU time breakdown: compute vs.\ transfer}
\label{fig:tsp_gpu_breakdown}
\end{subfigure}
\end{figure}

The speedup difference (1.34$\times$ vs.\ 1.13$\times$) comes from the CPU side: the original GA class runs about 32\% slower (664.2\,ms vs.\ 504.5\,ms) because of per-generation overhead from logging, fitness history tracking, and object allocation. The library uses the same GA logic but with less overhead. GPU times are similar across frameworks (Figure~\ref{fig:tsp_wall_time}), so this is not a GPU regression. Figure~\ref{fig:tsp_gpu_breakdown} shows transfer overhead at about one quarter of GPU time in both cases.

The fitness $p$-values match exactly ($p = 0.0072$): the GPU produces better tours than the CPU in both frameworks. Effect sizes are similar ($-435.7$ in the library vs.\ $-419.2$ in the original), suggesting the GPU mutation kernel explores the search space more broadly. As with XOR, warm-up reduces GPU standard deviation (184.8\,ms $\to$ 75.3\,ms) and narrows the confidence interval ([1.19$\times$, 1.46$\times$] $\to$ [1.08$\times$, 1.19$\times$]).

\begin{figure}[H]
\centering
\begin{subfigure}{0.49\textwidth}
\includegraphics[width=\textwidth]{diagrams/sec4/comparison_throughput_tsp.png}
\caption{Fitness evaluation throughput}
\label{fig:tsp_throughput}
\end{subfigure}
\hfill
\begin{subfigure}{0.49\textwidth}
\includegraphics[width=\textwidth]{diagrams/sec4/comparison_run_by_run_tsp.png}
\caption{Run-by-run wall times (50 paired runs)}
\label{fig:tsp_run_by_run}
\end{subfigure}
\end{figure}

Figure~\ref{fig:tsp_throughput} shows a 16\% GPU throughput difference (186K vs.\ 157K evals/sec), which is run-to-run GPU variance. Figure~\ref{fig:tsp_run_by_run} shows stable per-run times in both frameworks (warm-up removes the first-run JIT spike). The CPU times have a consistent offset, with the original framework always slower, which supports the overhead explanation.

To see where GPU wall time goes, the library was profiled with cProfile and NVIDIA Nsight Systems. Table~\ref{tab:gpu_profile} breaks down a single post-warm-up run for each problem.


\subsubsection*{GPU Execution Profiling}

\begin{table}[H]
\centering
\caption{GPU wall time breakdown by component (single post-warm-up run, library)}
\label{tab:gpu_profile}
\begin{tabular}{lcccc}
\toprule
 & \multicolumn{2}{c}{\textbf{XOR (640\,ms)}} & \multicolumn{2}{c}{\textbf{TSP (669\,ms)}} \\
\cmidrule(lr){2-3} \cmidrule(lr){4-5}
\textbf{Component} & ms & \% & ms & \% \\
\midrule
Kernel dispatch + GPU work & 377 & 58.9 & 47 & 7.0 \\
Host$\to$device transfer & 132 & 20.6 & 89 & 13.3 \\
Device$\to$host transfer & 63 & 9.8 & 37 & 5.5 \\
CPU-side GA logic & 5 & 0.8 & 303 & 45.3 \\
Python / Numba overhead & 63 & 9.8 & 193 & 28.9 \\
\bottomrule
\end{tabular}
\end{table}

The two problems have opposite bottlenecks. XOR is \emph{kernel-dominated}: GPU dispatch takes 59\% of wall time, transfers another 30\%. TSP is \emph{CPU-dominated}: crossover and tournament selection run on the host and take 45\%, while GPU kernels account for only 7\%. Nsight Systems shows that raw CUDA API costs are small (kernel launch: 12--31\,$\mu$s, DMA transfers: 37--138\,$\mu$s). The gap between these and the cProfile numbers comes from Numba's Python-side dispatch: each kernel call adds roughly 150--200\,$\mu$s of context checks and type inference. This overhead is part of Numba itself and applies equally to both frameworks.


\subsubsection*{Common Findings}

\begin{figure}[H]
\centering
\includegraphics[width=0.75\textwidth]{diagrams/sec4/comparison_speedup.png}
\caption{GPU speedup with 95\% bootstrap CI across both benchmarks}
\label{fig:comparison_speedup}
\end{figure}

Figure~\ref{fig:comparison_speedup} shows GPU speedup across both benchmarks. All four experiments (two frameworks, two problems) show statistically significant acceleration, with every confidence interval above the breakeven line. The speedup differences between frameworks (1.51$\times$ vs.\ 1.39$\times$ for XOR; 1.34$\times$ vs.\ 1.13$\times$ for TSP) come from CPU-side overhead, not GPU regressions, since GPU wall times are similar within each benchmark.

Transfer overhead is consistent across frameworks: 29.2\% vs.\ 31.5\% for XOR, 23.8\% vs.\ 25.1\% for TSP. This shows that pinned-memory transfers and kernel execution work the same way in both implementations.

Fitness correctness holds in all experiments. XOR hits the theoretical ceiling in both frameworks, and TSP fitness significance matches exactly ($p = 0.0072$), with GPU tours achieving lower path costs. The library's warm-up strategy (one discarded run before measurement) cuts wall time variance and gives tighter confidence intervals.

\end{document}
